\section{Methodology}
\label{sec:method}

We compare three paradigms for accessing stylistic distributions in language models:
(1)~direct generation from the base model via completion-style prompting,
(2)~instruction-based generation from the aligned model, and
(3)~logit-space distribution arithmetic that interpolates between the two.

\subsection{Models and Setup}
\label{sec:models}

We use \qwenbase as the base model and \qweninst as the aligned model, both with
3 billion parameters. Each model is loaded in float16 on a separate NVIDIA RTX
A6000 GPU (49\,GB), enabling simultaneous inference for distribution arithmetic
with KV caching.

\subsection{Style Tasks}
\label{sec:styles}

We select seven target styles spanning literary and functional registers:
five authorial styles (Hemingway, Shakespeare, Poe, Austen, Tolkien) and two
functional styles (academic writing, film noir narration). These styles were
chosen to span a range of distinctiveness---from highly ornate and recognizable
(Shakespeare, Poe) to understated and subtle (Hemingway)---allowing us to test
whether some styles are easier to invoke than others.

Each style is applied to three topics: (1)~``A person arriving at an old house
for the first time,'' (2)~``The experience of watching a sunset over the ocean,''
and (3)~``A conversation between two strangers on a train.'' This yields
$7 \times 3 = 21$ style--topic tasks.

\subsection{Generation Conditions}
\label{sec:conditions}

For each style--topic task, we generate text under five conditions:

\para{Base model ($\interpcoeff = 0$).}
The base model receives a completion-style prompt describing the target style and
topic. We generate 3 samples per task.

\para{Aligned model ($\interpcoeff = 1$).}
The aligned model receives a chat-formatted instruction requesting text in the
target style. We generate 3 samples per task.

\para{Distribution arithmetic ($\interpcoeff \in \{0.25, 0.5, 0.75\}$).}
At each decoding step, we compute logits from both models on the same prefix and
combine them linearly:
\begin{equation}
    \ell_{\text{mix}} = (1 - \interpcoeff) \cdot \ell_{\text{base}} + \interpcoeff \cdot \ell_{\text{align}}
    \label{eq:interp}
\end{equation}
where $\ell_{\text{base}}$ and $\ell_{\text{align}}$ are the raw logit vectors
from the base and aligned models, respectively, and $\interpcoeff$ controls the
mixture weight. Sampling then proceeds from $\text{softmax}(\ell_{\text{mix}})$.
We generate 2 samples per task at each interpolation level.

All conditions use temperature $\tau = 0.8$, top-$p = 0.95$, maximum length of
128 tokens, and random seed 42. This yields $21 \times (3 + 3 + 2 \times 3) = 231$
total text samples.

\subsection{Evaluation}
\label{sec:eval}

\para{LLM-as-judge.}
We use \gptjudge to evaluate each sample on two dimensions: \emph{style fidelity}
(1--10 scale: how well the output captures the target style's distinctive voice,
vocabulary, and sentence structure) and \emph{text quality} (1--10 scale: coherence
and writing quality regardless of style). This required 210 API calls across all
condition--task pairs.

\para{Lexical diversity.}
We compute Distinct-$N$ (type-token ratio for $N$-grams, $N \in \{1,2,3\}$) to
measure vocabulary diversity within each sample. Higher values indicate more
diverse word choices.

\para{Intra-condition diversity.}
We compute Self-BLEU-4 (4-gram overlap between samples within each condition--style
pair). Lower values indicate that samples within the same condition are more
diverse from each other.

\para{Style differentiation.}
We apply one-way ANOVA on style fidelity scores across the seven target styles
within each condition. A significant $F$-statistic indicates that the model
produces outputs that are distinguishable across styles---\ie it is not generating
the same generic text regardless of the style instruction.

\para{Statistical tests.}
We use Mann-Whitney $U$ tests for pairwise comparisons between conditions and
report $p$-values with standard significance thresholds ($^{*}p < 0.05$,
$^{**}p < 0.01$, $^{***}p < 0.001$).
